%\title{Building a Multilayer Perceptron (MLP) Network from scratch}
%\maketitle
%\chapter*{Definition of an Artificial Neural Network (ANN)}
%A neural network is an approximator whose complexity arises through the intertwining of simple components, not through few complicated ones. The connection of the simple components, called \textit{nodes}, is systematic which can be exploited during calculations.\par
%ANNs are often described as a model of the human brain which is neither correct nor helpful, so that analogy will never be used here. A text based definition can only go so far.
\chapter{Multilayer Perceptron Network}
The simplest variant of an ANN is the MLP. Its applications are limited, because a pure MLP does not handle gigantic inputs very well like convolutional nets (CNN) and it has no inherent way of representing timeseries like recurrent nets and long-short-term -memory nets (RNN, LSTM). Although application of MLP nets have been deployed successfully as policy agents in state space controllers learned by reinforcement methods. Learning how MLPs function is the foundation on what to build next.\par
An MLP network consists of $L$ layers. The network layout $l$  is:
\begin{align}
	l = \begin{bmatrix}N_1&\cdots&N_{L}\end{bmatrix}
\end{align}
$N_l$ is the number of nodes in the $l-th$ layer.
\section{Transition between layers}
In a \textit{dense} network, each node of layer $l-1$ is connected to each node of layer $l$. A node is primarily characterized by the kind of activation function it uses.
\begin{align}
	\label{eq:activations}
	a_i^{l} =f\bra{z_i^{l}}
\end{align}
$a_i^{l}$ is the activation of the $i$-th node in the $l$-th layer. $z_i^{l}$ is the pre-activation input of the $i$-th node in the $l$-th layer. $z_i^{l}$ is a linear combination of the activations of the previous layer.
\begin{align}
	\label{eq:preactivationinput}
	z_i^{l}=\sum\limits_{j=0}^{N_{l-1}-1}w_{ij}^{l-1}\cdot a_j^{l-1}+b_i^{l-1}
\end{align}

Note that the bias term $b_i^{l-1}$ belongs to the $l-1$-th transition or the $l-1$-th layer. To add it to the $l$-th layer  is just confusing and makes index housekeeping harder.
Because it is used in textbooks, from the sum-component form, a matrix form can be derived. First, the sum is made into the multiplication of a row vector with a column vector.
\begin{align}
	\label{eq:vectorproduct1}
	z_i^{l}=\boldsymbol{w}_i^{l-1,\top}\cdot \boldsymbol{a}^{l-1}+b_i^{l-1}
\end{align}
Expanding this for every node in the $l$-th layer, we get:
\begin{align}
	\begin{bmatrix}z_0^{l}\\\vdots\\z_{N_l-1}^{l}\end{bmatrix}&=\begin{bmatrix}\boldsymbol{w}_0^{l-1,\top}\cdot \boldsymbol{a}^{l-1}\\\vdots\\\boldsymbol{w}_{N_l-1}^{l-1,\top}\cdot \boldsymbol{a}^{l-1}\end{bmatrix}+\begin{bmatrix}b_{N_l-1}^{l-1}\\\vdots\\ b_{N_l-1}^{l-1}\end{bmatrix}=\begin{bmatrix}\boldsymbol{w}_0^{l-1,\top}\\\vdots\\\boldsymbol{w}_{N_l-1}^{l-1,\top}\end{bmatrix} \cdot\boldsymbol{a}^{l-1}+\boldsymbol{b}^{l-1}\\&=\boldsymbol{W}^{l-1} \cdot\boldsymbol{a}^{l-1}+\boldsymbol{b}^{l-1}
\end{align}
\section{Computing the forward pass efficiently}
\eqref{eq:activations} and \eqref{eq:preactivationinput} actually contain all information to understand the forward pass. But the goal is to implement it efficiently. That means storing a matrix for every layer is not a good idea. The first reason is that in general the dimensions vary, therefore no stack allocations can be used to store the array of arrays on the stack, meaning there instantly is a double indirection for every read. This could be mitigated by using a linar algebra library but the task is not to use libraries but to think about an efficient layout. Also, a flattened array with constant time access still is faster than singular flattened matrices. For small to medium size networks it is even plausible to assume that the contiguous parameter array will land in a cache.
\subsection{Bias as constant input}
The bias term in \eqref{eq:preactivationinput} might lead to a better intuition for what is happening but not for an easier flattening of the parameter vector. We rewrite the equation
\begin{align}
	\label{eq:constantbias}
	z_i^{l}=\sum\limits_{j=0}^{N_{l-1}-1}w_{ij}^{l-1}\cdot a_j^{l-1}
\end{align}
Later, we force $a_{N_{l-1}}^{l-1}$ to be constant $1$. This makes \eqref{eq:constantbias} truly linear. The trick now is, just to add 1 to every node count except the last. This leads mathematically to a bias-free network which is easier to handle.
\subsection{Flattening of the parameter vector}
For a MLP net with a layout
\begin{align}
	l = \begin{bmatrix}
		N_1&N_2&\cdots&N_L
	\end{bmatrix}
\end{align}
the total number of parameters $	N_\Theta$ can be calculated as 
\begin{align}
	N_\Theta = \sum\limits_{i=1}^{\mathrm{dim}\bra{l}-1}l_{i+1}\cdot\bra{l_{i}+1}
\end{align}
We also can calculate the vector-offset of each transition with a similar formula
\begin{align}
	o_{i+1}=o_i+l_{i+1}\cdot\bra{l_{i}+1}
\end{align}
with $o_0=0$.
\section{Deriving Backpropagation}
Backpropagation is the gold standard algorithm in training neural networks. It is easy in concept, messy in derivation and then again easy to implement. \par
\enquote{Training a neural network} is in almost all cases exactly the same as solving a regression problem. From a certain point of view, it is nothing more than calculating the parameters of a regression line. Backpropagation is the algorithm that exploits the structure of an MLP heavily to make the nonlinear regression problem feasibly solvable.
\subsection{Loss}
We use the sum of squares of output errors as a loss function to make the analogy to any old regression even more pronounced. Other loss functions are possible but they all use the derivatives of output neurons w.r.t. the weights at some point. The loss is defined as:
\begin{align}
	\label{eq:loss}
	L&=\sum\limits_{k=0}^{N-1}\sum\limits_{i=0}^{N_L-1}\bra{a^L_i\bra{k}-y_i\bra{k}}^2
\end{align}
with $a^L_i\bra{k}-y_i\bra{k}=e_{ik}$
\subsection{Loss gradient}
If the partial derivative w.r.t. some weight of some layer is taken, the recursion tree will explode and it is not feasible for any net of any non-toy size net gradient to be calculated that way. The trick of backpropagation is to use the chain rule layer-wise. The trick has three separate components.
\begin{itemize}[]\item Expressing the Loss gradient w.r.t. weights in terms of an so-called \enquote{error signal}
	\item Calculating error signals not directly but with a (back)-propagation equation
	\item Only calculating the first error signal
\end{itemize}


A node's sample error signal is defined as the partial derivative of the loss w.r.t. to the pre-activations \eqref{eq:preactivationinput}:
\begin{align}
	\delta^p_i\bra{k}=\frac{\partial L}{\partial z_i^p(k)}
\end{align}
\subsubsection{Update equation}
An error signal arises naturally when targeting the gradient components of the loss w.r.t. the weights of the last transition - which are what is actually needed for any kind of gradient descent - and apply the chain rule (including total differentiation) to pull the definition of the error signal out.
\begin{align}
	\label{eq:chainrule-error-signal}
	\ptl{L}{w_{qr}^{p-1}}=\sum\limits_{k=0}\underset{\delta_{bk}^p}{\underbrace{\ptl{L}{z_{bk}^p}}}\cdot\ptl{z_{bk}^p}{w_{qr}^{p-1}}
\end{align}

Using \eqref{eq:constantbias}, the actual content of the second factor can be calculated.
\begin{align}
	\ptl{z_{bk}^p}{w_{qr}^{p-1}}&=\ptl{}{w_{qr}^{p-1}}\bra{\sum\limits_{j=0}^{N_{p-1}-1}w_{bj}^{p-1}\cdot a_{jk}^{p-1}}\\
	&=\sum\limits_{j=0}^{N_{p-1}-1}\ptl{w_{bj}^{p-1}}{w_{qr}^{p-1}}\cdot a_{jk}^{p-1}\\
	&=\sum\limits_{j=0}^{N_{p-1}-1}\delta^b_q\delta^j_ra_{jk}^{p-1}\\
	&=\delta^b_qa_{rk}^{p-1}
\end{align}
It might be obvious for some, but it is helpful to not assume the \enquote{correct} index in the derivative itself. There is only a nonzero derivative obtained if we differentiate the node w.r.t. the correct weights. This behavior has to arise from generalized formulas. The neuron $b$ is the output we are testing for and the $q$ indices are the destination paths to the layer. Therefore the weights with the first index $q$ don't contribute to neuron $b$.\par
With $b\neq q$ \eqref{eq:chainrule-error-signal} always equals zero, therefore we can comfortably write:
\begin{align}
	\ptl{L}{w_{br}^{p-1}}&=\sum\limits_{k=0}^{N-1}\underset{\delta_{bk}^p}{\underbrace{\ptl{L}{z_{bk}^p}}}\cdot\ptl{z_{bk}^p
	}{w_{br}^{p-1}}\\
	\label{eq:update-equation}
	&=\sum\limits_{k=0}^{N-1}\delta^p_{bk}\cdot a_{rk}^{p-1}
\end{align}
Note the important result here: \eqref{eq:update-equation} was derived with the last layer in mind, but this is the general equation for obtaining gradient components in any layer.
\subsubsection{Backpropagation equation}
We can focus on a single sample's error signal now. The sample index $k$ will be omitted until the very end, but every activation, pre-activation and error signal are sample-wise! Since want to derive a back propagation scheme, where we start at the end index, it is sound to target the next lower index. By definition is:
\begin{align}
	\delta_b^{p-1}=\ptl{L}{z_b^{p-1}}
\end{align}
Application of the chain rule yields:
\begin{align}
	\label{eq:backprop-full}
	\delta_b^{p-1}=\ptl{L}{a_b^{p-1}}\ptl{a_b^{p-1}}{z_b^{p-1}}
\end{align}
The second factor can be handled with \eqref{eq:activations}. It just becomes $f^\prime\bra{z_b^{p-1}}:=g_b^{p-1}$ which will be shortened to save inices and brackets. The next step is the actual step where the backpropagation takes place. Using the chain rule (including the total differentiation), the loss derivative w.r.t. the activations of layer $p-1$ can be expressed in terms of the node inputs of layer $p$. At some point in the network, this relationship exists. The chain rule is applied selectively here. Since the activation of neuron $b$ of layer $p-1$ contribute to all pre-activation inputs  of layer $p$, and all neurons in layer $p$ contribute to the loss at some point, the contributions of all neurons of layer $b$ have to be collected.
\begin{align}
	\ptl{L}{a_b^{p-1}}&=\sum\limits_{v=0}^{N_p-1}\ptl{L}{z_v^{p}}\ptl{z_v^{p}}{a_b^{p-1}}\\
	\label{eq:backprop-part}
	&=\sum\limits_{v=0}^{N_p-1}\delta_v^p\ptl{z_v^{p}}{a_b^{p-1}}
\end{align}
Using \eqref{eq:constantbias} the second term can be calculated:
\begin{align}
	\label{eq:backprop-part2}
	\ptl{z_v^{p}}{a_b^{p-1}}&=\ptl{}{a_b^{p-1}}\bra{\sum\limits_{j=0}^{N_{p-1}-1}w_{vj}^{p-1}\cdot a_j^{p-1}}=w_{vb}^{p-1}
\end{align}
Plugging \eqref{eq:backprop-part2} in \eqref{eq:backprop-part} and then into \eqref{eq:backprop-full} yields the full backpropagation equation in coordinate form with sample index $k$ restored.
\begin{align}
	\label{eq:backprop-final}
	\delta_{bk}^{p-1}=\bra{\sum\limits_{v=0}^{N_p-1}\delta_{vk}^pw_{vb}^{p-1}}g_{bk}^{p-1}
\end{align}
\subsubsection{Deriving the first error signal}
To derive the first sample error signal we simply take the derivative of the loss \eqref{eq:loss} w.r.t. the $b$-th pre-activation input to the $L$-th layer of the $k$-th sample and apply the chain rule analogous to \eqref{eq:backprop-full}.
\begin{align}
	\delta_{bk}^L=\ptl{L}{z^L_{bk}}=\ptl{L}{a_{bk}^{L}}\ptl{a_{bk}^{L}}{z_{bk}^{L}}=\ptl{L}{a_{bk}^{L}}g_{bk}^L
\end{align}
Since $L$ is the output layer the activations of layer $L$ actually exist in the Loss function directly such that there is no more layer traversal necessary.
\begin{align}
	\label{eq:first-signal}
	\delta_{bk}^L=e_{bk}\cdot g_{bk}^L
\end{align}
\subsection{From Coordinates to Contiguous Memory}
While the coordinate-based derivation provides the necessary rigor to understand the flow of error through the network, it is a poor blueprint for a high-performance implementation. In languages like Rust, treating layers as a collection of separate heap-allocated objects introduces significant overhead through fragmentation, bounds-checking, and the destruction of CPU cache locality.

To move from abstraction to a "bare-metal" implementation, we perform two transformations:

\begin{enumerate}
	\item \textbf{Vector-Matrix Notation:} We group coordinate-wise operations into linear algebra primitives. The forward pass becomes a series of affine transformations, and backpropagation becomes a sequence of transposed matrix-vector multiplications.
	\item \textbf{Linear Flattening:} To achieve maximum throughput, we treat the entire network as a \textbf{single, contiguous block of memory}. By mapping topological indices $(p, b, r)$ to linear offsets within one \texttt{Vec<f32>} (or a static array for microcontrollers), we eliminate the need for general-purpose libraries.
\end{enumerate}

This approach transforms the training loop into a predictable stream of Multiply-Accumulate (MAC) operations, allowing the CPU prefetcher to operate at peak efficiency and simplifying Gradient Descent into a single, unified vector subtraction.
\subsection{Loss gradient (vectorized)}
\subsubsection{Update equation (vectorized)}
Since we exploit the layer-wise nature, we symbolically write the following derivative based on the component representation \eqref{eq:update-equation}.
\begin{align}
	\label{eq:loss-gradient-vectorized}
	\ptl{L}{\boldsymbol{W}^{p-1}}&=\sum\limits_{k=0}^{N-1}\begin{bmatrix}
		\delta_{0k}^p\cdot a_{0k}^{p-1}&\cdots&\delta_{0k}^p\cdot a_{N_{p-1}-1,k}^{p-1}\\
		\vdots&\ddots&\vdots\\
		\delta_{N_{p}-1,k}^p\cdot a_{0k}^{p-1}&\cdots&\delta_{N_{p}-1,k}^p\cdot a_{N_{p-1}-1,k}^{p-1}\\
	\end{bmatrix}=\sum\limits_{k=0}^{N-1}\boldsymbol{\delta}^p_k\otimes\boldsymbol{a}^{p-1}_k
\end{align}

\eqref{eq:loss-gradient-vectorized} is the notation many textbooks use without a rigorous handling of the indices. It is great for transporting the idea for a semi-technical audience but is insufficient for people who want to understand the structure completely. Hardware does not know about matrices, only about linear addresses and offsets. If we bake that step into the algorithm directly, we save layers of indirection. What we actually want is a flattened gradient, not a matrix of partial derivatives. For that, we observe that the matrix can be written as
\begin{align}
	&\begin{bmatrix}
		\delta_{0k}^p\cdot a_{0k}^{p-1}&\cdots&\delta_{0k}^p\cdot a_{N_{p-1}-1,k}^{p-1}\\
		\vdots&\ddots&\vdots\\
		\delta_{N_{p}-1,k}^p\cdot a_{0k}^{p-1}&\cdots&\delta_{N_{p}-1,k}^p\cdot a_{N_{p-1}-1,k}^{p-1}
	\end{bmatrix}\\
	&\Rightarrow\begin{bmatrix}\boldsymbol{\delta}^p_k\cdot a_{0k}^{p-1} & \cdots &\boldsymbol{\delta}^p_k\cdot a_{N_{p-1}-1,k}^{p-1} \end{bmatrix}\\&=\begin{bmatrix}\boldsymbol{\delta}^p_k\cdot a_{0k}^{p-1} \\ \vdots \\ \boldsymbol{\delta}^p_k\cdot a_{N_{p-1}-1,k}^{p-1} \end{bmatrix}=	\ptl{L}{\boldsymbol{W}^{p-1}_\mathrm{flat}}
\end{align}
The last step is a rearangement of the Matrix from a being $\in \mathbb{R}^{N_p\times N_{p-1}}$ to $\in \mathbb{R}^{N_p\cdot N_{p-1}\times 1}$. This gradient can be stored in an contiguous array for constant time access once its loaded.
\subsubsection{Backpropagation (vectorized)}
This can be done in two steps. From \eqref{eq:backprop-final} we first identify that the summation can be represented as another dot product
\begin{align}
	\label{eq:backprop-vectorized}
	\delta_{bk}^{p-1}=\begin{bmatrix}\delta_{0k}^{p}&\cdots&\delta_{N_{p}-1,k}^{p}\end{bmatrix}\begin{bmatrix}
		w_{0b}^{p-1}\\ \vdots \\ w_{N_{p}-1,b}^{p-1}
	\end{bmatrix} g_{bk}^{p-1}=\boldsymbol{\delta}_k^{p,\top}\cdot \boldsymbol{w}_b^{p-1}\cdot  g_{bk}^{p-1}
\end{align}

$\boldsymbol{w}_n^{p-1}$ is the $b$-th column of the weight matrix of the $p-1$-th transition. Now, we can construct the left-hand side to be the full error signal vector.

\begin{align}
	\label{eq:backprop-vectorized2}
	\delta_{bk}^{p-1}&=\boldsymbol{\delta}_k^{p,\top}\cdot \boldsymbol{w}_b^{p-1}\cdot  g_{bk}^{p-1}\\
	&\Downarrow\nonumber\\ 
	\begin{bmatrix}
		\delta_{0k}^{p-1}&\cdots&\delta_{N_{p-1}-1,k}^{p-1}
	\end{bmatrix}&=	\begin{bmatrix}
		\boldsymbol{\delta}_k^{p,\top}\cdot \boldsymbol{w}_0^{p-1}\cdot  g_{0k}^{p}&\cdots&\boldsymbol{\delta}_k^{p-1,\top}\cdot \boldsymbol{w}_{N_{p-1}-1}^{p-1}\cdot  g_{N_{p-1}-1,k}^{p-1}\\
	\end{bmatrix}\\
	&\Downarrow\nonumber\\ 
	\begin{bmatrix}
		\delta_{0k}^{p-1}\\\vdots\\\delta_{N_{p-1}-1,k}^{p-1}
	\end{bmatrix}&=	\begin{bmatrix}
		\boldsymbol{\delta}_k^{p,\top}\cdot \boldsymbol{w}_0^{p-1}\cdot  g_{0k}^{p-1}\\\vdots\\\boldsymbol{\delta}_k^{p,\top}\cdot \boldsymbol{w}_{N_{p-1}-1}^{p-1}\cdot  g_{N_{p-1}-1,k}^{p-1}\\
	\end{bmatrix}\\
	\boldsymbol{\delta}_k^{p-1}&=\begin{bmatrix}
		\boldsymbol{w}_0^{p-1,\top}\cdot \boldsymbol{\delta}_k^{p}\cdot  g_{0k}^{p-1}\\\vdots\\ \boldsymbol{w}_{N_{p-1}-1,\top}^{p-1}\cdot\boldsymbol{\delta}_k^{p}\cdot  g_{N_{p-1}-1,k}^{p-1}\\
	\end{bmatrix}
\end{align}
The last equation is rewritten with the hadamard product.
\begin{align}
	\boldsymbol{\delta}_k^{p-1}&=\begin{bmatrix}
		\boldsymbol{w}_0^{p-1,\top}\cdot \boldsymbol{\delta}_k^{p}\\\vdots\\ \boldsymbol{w}_{N_{p-1}-1,\top}^{p-1}\cdot\boldsymbol{\delta}_k^{p}\\
	\end{bmatrix}\odot\begin{bmatrix}
		g_{0k}^{p-1}\\ \vdots \\
		g_{N_{p-1}-1,k}^{p-1}
	\end{bmatrix}\\
	&=\left[\underset{\boldsymbol{W}^{p-1,\top}}{\underbrace{\begin{bmatrix}
				\boldsymbol{w}_0^{p-1,\top}\\\vdots\\ \boldsymbol{w}_{N_{p-1}-1,\top}^{p-1}\\
	\end{bmatrix}}}\cdot\boldsymbol{\delta}_k^{p}\right]\odot\boldsymbol{g}_k^{p-1}\\
	\label{eq:error_signal_update}
	&=\boldsymbol{W}^{p-1,\top}\cdot\boldsymbol{\delta}_k^{p}\odot\boldsymbol{g}_k^{p-1}
\end{align}
This is still a matrix multiplication that needs some adjustment to be closer to the actual implementation. But there is no direct representation that codifies that adjustment. The implementation of the parameter vector access in the backpropagation step will be addressed in \ref{sec:backpropagation-impl}
\subsubsection{The first error signal (vectorized)}
\eqref{eq:first-signal} can be directly written with the hadamard product:
\begin{align}
	\underset{\boldsymbol{\delta}_k^L}{\underbrace{	\begin{bmatrix}
				\delta_{0k}^L\\\vdots\\\delta^L_{N_{L}-1,k}
	\end{bmatrix}}}&=\underset{\boldsymbol{e}_k}{\underbrace{\begin{bmatrix}
				e_{0k}\\\vdots\\e_{N_{L}-1,k}
	\end{bmatrix}}}\odot\underset{\boldsymbol{g}_k^L}{\underbrace{	\begin{bmatrix}
				g{0k}^L\\\vdots\\g^L_{N_{L}-1,k}
	\end{bmatrix}}}\\
	\boldsymbol{\delta}_k^L&=\boldsymbol{e}_k\odot \boldsymbol{g}_k^L
\end{align}
\subsection{Backpropagation algorithm}
With all pieces in place, the procedure can be defined. The situation is that there are in total $N_\mathrm{total}$ samples available for the network. Those samples have been divided into a number of batches of size $N$. The gradient will be calculated batch-wise.
\begin{algorithm}[H]
	\SetAlgoLined
	\DontPrintSemicolon
	\caption{Mini-batch Gradient Descent with Epoch Cycling}
	
	\KwIn{Training set $\mathcal{D}$, Batch size $N$, Learning rate $\eta$}
	\BlankLine
	
	\While{termination criterion not met}{
		\textbf{Shuffle} dataset $\mathcal{D}$\;
		\ForEach{batch $\mathcal{B} \subset \mathcal{D}$ of size $N$}{
			\tcc{Parallelization over samples $k \in \mathcal{B}$}
			\ForEach{sample $k$}{
				1. \textbf{Forward Pass}: Compute and store activations $a^l_i$\;
				2. \textbf{Backward Pass}: Compute and store error signals $\delta^l_i$\;
				3. \textbf{Local Gradient}: Calculate weight update contribution $\Delta w_k$\;
			}
			4. \textbf{Aggregation}: Sum contributions $\Delta W = \sum_{k=1}^{N} \Delta w_k$\;
			5. \textbf{Update}: Adjust weights $w \leftarrow w - \eta \Delta W$\;
		}
		\textbf{Check Termination}: Evaluate loss/convergence to decide if another epoch is needed\;
	}
\end{algorithm}


















